% !TEX root = exercices_s2.tex
% ============================================================================
% HEC Montréal MBA -- ECON50803: Macroeconomic Environment
% Shared Preamble for Exercise Handouts
% ============================================================================

% --- Document class ---------------------------------------------------------
\documentclass[11pt, letterpaper]{article}
\usepackage[margin=2.5cm]{geometry}

% --- Fonts ------------------------------------------------------------------
\usepackage[sfdefault,light]{FiraSans}
\usepackage{FiraMono}
\usepackage[T1]{fontenc}

% --- Colors (same as slides) ------------------------------------------------
\usepackage{xcolor}
\definecolor{HECnavy}{HTML}{002855}
\definecolor{HECgreen}{HTML}{26d07c}
\definecolor{HECcoral}{HTML}{ff585d}
\definecolor{HECtext}{HTML}{2D3748}
\definecolor{HECgray}{HTML}{94A3B8}
\definecolor{HEClightgray}{HTML}{F1F5F9}
\definecolor{HEClightblue}{HTML}{E8EDF3}

% --- Packages ---------------------------------------------------------------
\usepackage{amsmath, graphicx, tikz, booktabs, tabularx, hyperref}
\usepackage[inline]{enumitem}
\usepackage{etoolbox}
\usepackage{comment}
\usepackage{needspace}
\usepackage{tcolorbox}
\tcbuselibrary{skins, breakable}
\usepackage{titlesec}
\usepackage{fancyhdr}
\usepackage{lastpage}

% --- Paths ------------------------------------------------------------------
\graphicspath{{../Slides/Figures/}}

% --- Hyperlinks -------------------------------------------------------------
\hypersetup{
   colorlinks,
   linkcolor=HECnavy,
   urlcolor=HECnavy,
   citecolor=HECtext
}

% --- Answer toggle ----------------------------------------------------------
\newbool{answers}
\ifdefined\showanswers
   \booltrue{answers}
\else
   \boolfalse{answers}
\fi

% Answer environment: shown as green box when answers=true, excluded otherwise
\ifbool{answers}{%
   \newtcolorbox{reponse}[1][Réponse]{%
      colback=HECgreen!5,
      colframe=HECgreen,
      boxrule=0pt,
      leftrule=3pt,
      arc=4pt,
      title={\textcolor{HECgreen}{\bfseries #1}},
      coltitle=HECgreen,
      fonttitle=\bfseries,
      attach title to upper={\\[4pt]},
      left=8pt, right=8pt, top=4pt, bottom=4pt,
      before={\par\nopagebreak\medskip\nopagebreak}
   }%
}{%
   \excludecomment{reponse}%
}

% --- Custom box environments (adapted from slides) --------------------------

% Key Insight (green accent)
\newtcolorbox{keyinsight}[1][Point clé]{%
   colback=HECgreen!5,
   colframe=HECgreen,
   boxrule=0pt,
   leftrule=3pt,
   arc=4pt,
   title={\textcolor{HECgreen}{\bfseries #1}},
   coltitle=HECgreen,
   fonttitle=\bfseries,
   attach title to upper={\\[4pt]},
   left=8pt, right=8pt, top=4pt, bottom=4pt,
   breakable
}

% Business Implication (navy accent)
\newtcolorbox{bizimplication}[1][Implication d'affaires]{%
   colback=HECnavy!5,
   colframe=HECnavy,
   boxrule=0pt,
   leftrule=3pt,
   arc=4pt,
   title={\textcolor{HECnavy}{\bfseries #1}},
   coltitle=HECnavy,
   fonttitle=\bfseries,
   attach title to upper={\\[4pt]},
   left=8pt, right=8pt, top=4pt, bottom=4pt,
   breakable
}

% Warning / Important (coral accent)
\newtcolorbox{warning}[1][Important]{%
   colback=HECcoral!5,
   colframe=HECcoral,
   boxrule=0pt,
   leftrule=3pt,
   arc=4pt,
   title={\textcolor{HECcoral}{\bfseries #1}},
   coltitle=HECcoral,
   fonttitle=\bfseries,
   attach title to upper={\\[4pt]},
   left=8pt, right=8pt, top=4pt, bottom=4pt,
   breakable
}

% Definition (gray background)
\newtcolorbox{defbox}[1][Définition]{%
   colback=HEClightgray,
   colframe=HEClightgray,
   boxrule=0pt,
   leftrule=3pt,
   arc=4pt,
   left=8pt, right=8pt, top=4pt, bottom=4pt,
   title={\textcolor{HECnavy}{\bfseries #1}},
   coltitle=HECnavy,
   fonttitle=\bfseries,
   attach title to upper={\\[4pt]},
   breakable
}

% Exercise (coral accent)
\newtcolorbox{exercisebox}[1][Exercice]{%
   colback=HECcoral!5,
   colframe=HECcoral,
   boxrule=0pt,
   leftrule=3pt,
   arc=4pt,
   title={\textcolor{HECcoral}{\bfseries #1}},
   coltitle=HECcoral,
   fonttitle=\bfseries,
   attach title to upper={\\[4pt]},
   left=8pt, right=8pt, top=4pt, bottom=4pt,
   breakable
}

% --- Section formatting (titlesec) ------------------------------------------
\titleformat{\section}
   {\clearpage\Large\bfseries\color{HECnavy}}
   {\thesection.}{0.5em}{}
   [\vspace{2pt}{\color{HECgreen}\rule{2.5cm}{1.5pt}}]
\titlespacing*{\section}{0pt}{1.5em}{0.8em}

\titleformat{\subsection}
   {\large\bfseries\color{HECnavy}}
   {\thesubsection.}{0.5em}{}
\titlespacing*{\subsection}{0pt}{1em}{0.5em}

% --- Header/Footer (fancyhdr) ----------------------------------------------
\pagestyle{fancy}
\fancyhf{}
\renewcommand{\headrulewidth}{0pt}
\renewcommand{\footrulewidth}{0.4pt}
\renewcommand{\footrule}{\hfill{\color{HECgray}\rule{0.9\textwidth}{0.4pt}}\hfill\vspace{4pt}}
\fancyfoot[L]{\small\textcolor{HECgray}{ECON 50803\enspace\textbar\enspace Environnement macroéconomique}}
\fancyfoot[R]{\small\textcolor{HECgray}{\thepage\enspace/\enspace\pageref{LastPage}}}

% --- Enumitem defaults ------------------------------------------------------
\setlist[itemize]{leftmargin=1.5em, itemsep=3pt}
\setlist[enumerate]{leftmargin=1.5em, itemsep=3pt}

% --- MC question helpers ----------------------------------------------------
\newcounter{qcmcounter}
\newcommand{\qcm}[1]{%
   \stepcounter{qcmcounter}%
   \par\vspace{1.2em}%
   \needspace{12\baselineskip}%
   \noindent\textbf{\textcolor{HECnavy}{\theqcmcounter.}}\enspace #1%
   \vspace{0.3em}%
   \nopagebreak[4]%
}

\newcommand{\choix}[1]{%
   \nopagebreak[4]%
   \begin{enumerate}[label=\textcolor{HECgray}{(\alph*)}, leftmargin=2em, itemsep=1pt, topsep=2pt]
      #1
   \end{enumerate}
}

% --- Short answer counter ---------------------------------------------------
\newcounter{shortcounter}
\newcommand{\shortq}[1]{%
   \stepcounter{shortcounter}%
   \par\vspace{1.5em}%
   \needspace{4\baselineskip}%
   \noindent\textbf{\textcolor{HECnavy}{Question \theshortcounter.}}\enspace #1%
   \vspace{0.3em}%
}

% --- VFI (Vrai/Faux/Incertain) helper --------------------------------------
\newcommand{\vfi}[1]{%
   \stepcounter{shortcounter}%
   \par\vspace{1.5em}%
   \needspace{10\baselineskip}%
   \noindent\textbf{\textcolor{HECnavy}{Question \theshortcounter.}}\enspace\textit{Vrai, faux ou incertain?} #1%
   \vspace{0.3em}%
}

% --- Title page command -----------------------------------------------------
\newcommand{\exercisetitlepage}[2]{%
   % #1 = session number, #2 = session title
   \thispagestyle{empty}
   \begin{tikzpicture}[remember picture, overlay]
      % Navy sidebar
      \fill[HECnavy] (current page.north west) rectangle
         ([xshift=0.8cm]current page.south west);
      % Green accent line
      \fill[HECgreen] ([xshift=0.8cm]current page.north west) rectangle
         ([xshift=0.85cm]current page.south west);
   \end{tikzpicture}
   \vspace*{2cm}
   \hspace{0.8cm}%
   \begin{minipage}{0.85\textwidth}
      {\fontsize{28}{34}\selectfont\bfseries\textcolor{HECnavy}{Exercices --- Séance #1}}
      \\[14pt]
      {\Large\textcolor{HECnavy}{#2}}
      \\[14pt]
      {\color{HECgreen}\rule{3cm}{2pt}}
      \\[14pt]
      {\large\textcolor{HECgray}{ECON 50803 --- Environnement macroéconomique}}
      \\[8pt]
      {\large\textcolor{HECgray}{Jean-Félix Brouillette}}
      \\[20pt]
      \ifbool{answers}{%
         {\Large\bfseries\textcolor{HECgreen}{AVEC SOLUTIONS}}%
      }{%
         {\Large\bfseries\textcolor{HECgray}{SANS SOLUTIONS}}%
      }
   \end{minipage}
   \vfill
   \hspace{0.8cm}%
   {\small\textcolor{HECgray}{HEC Montr\'eal\enspace\textbar\enspace MBA}}
   \vspace{1cm}
   \newpage
}

% --- Mathematical commands (from slides) ------------------------------------
\newcommand{\pctchg}[1]{\%\Delta #1}


\begin{document}

\exercisetitlepage{2}{Croissance de long terme}

% ============================================================================
\section{Questions à choix multiples}
% ============================================================================

\setcounter{qcmcounter}{0}

\qcm{La fonction de production Cobb-Douglas est $Y = A \cdot K^{\alpha} \cdot L^{1-\alpha}$. Si on double \textbf{uniquement} le stock de capital ($K$), la production :}
\choix{
   \item Double exactement
   \item Plus que double
   \item Augmente, mais moins que double
   \item Reste inchangée
}

\begin{reponse}
\textbf{(c)} La fonction Cobb-Douglas présente des \textbf{rendements décroissants} par rapport à chaque facteur pris individuellement. Puisque $\alpha < 1$, doubler $K$ multiplie la production par $2^{\alpha} < 2$. Avec $\alpha = 0.3$, la production augmente d'un facteur $2^{0.3} \approx 1.23$, soit une hausse de 23\,\%, pas de 100\,\%.
\end{reponse}

\qcm{Si on double \textbf{simultanément} le capital ($K$) et le travail ($L$) dans la fonction Cobb-Douglas, la production :}
\choix{
   \item Double exactement
   \item Plus que double
   \item Augmente, mais moins que double
   \item Dépend de la valeur de $\alpha$
}

\begin{reponse}
\textbf{(a)} C'est la propriété des \textbf{rendements d'échelle constants} : $A \cdot (2K)^{\alpha} \cdot (2L)^{1-\alpha} = 2^{\alpha} \cdot 2^{1-\alpha} \cdot A K^{\alpha} L^{1-\alpha} = 2Y$. C'est l'argument de réplication : pour doubler la production d'un McDonald, il suffit d'en construire un autre identique de l'autre côté de la rue.
\end{reponse}

\qcm{La productivité totale des facteurs (PTF, notée $A$) :}
\choix{
   \item Se mesure directement dans les données nationales
   \item Est calculée comme un résidu après avoir soustrait les contributions de $K$ et $L$
   \item Correspond uniquement au progrès technologique
   \item Est toujours plus élevée dans les pays riches en ressources naturelles
}

\begin{reponse}
\textbf{(b)} La PTF est un \textbf{résidu} de la décomposition de Solow : $\pctchg{A} = \pctchg{Y} - \alpha \cdot \pctchg{K} - (1-\alpha) \cdot \pctchg{L}$. On n'observe pas $A$ directement; on la déduit après avoir mesuré $Y$, $K$, $L$ et $\alpha$. C'est une \guillemotleft~mesure de notre ignorance~\guillemotright. (c) est trop restrictif : $A$ capte la technologie, mais aussi le savoir-faire, les institutions, l'organisation, etc.
\end{reponse}

\qcm{Dans la fonction Cobb-Douglas $Y = A \cdot K^{\alpha} \cdot L^{1-\alpha}$ avec $\alpha = 0.3$, quelle part du revenu national va au travail?}
\choix{
   \item 30\,\%
   \item 50\,\%
   \item 70\,\%
   \item Cela dépend du niveau de $A$
}

\begin{reponse}
\textbf{(c)} Une propriété fondamentale de la Cobb-Douglas est que les parts du revenu national sont \textbf{constantes} et déterminées par les exposants : la part du capital est $\alpha = 0.3$ (30\,\%) et la part du travail est $1 - \alpha = 0.7$ (70\,\%). Ces parts ne dépendent ni de $A$, ni du niveau de $K$ ou $L$.
\end{reponse}

\qcm{Le cours utilise le ratio $K/Y$ plutôt que $K/L$ pour décomposer le PIB par habitant. Pourquoi?}
\choix{
   \item Parce que $K/Y$ est plus facile à mesurer
   \item Parce que $K/L$ attribue au capital de la croissance en réalité causée par la productivité ($A$)
   \item Parce que $K/Y$ est toujours constant dans le temps
   \item Parce que $K/L$ n'est pas disponible dans les données
}

\begin{reponse}
\textbf{(b)} Quand $A$ augmente, les entreprises investissent davantage (le capital suit la productivité). Utiliser $K/L$ attribuerait cette accumulation de capital \textbf{induite par} $A$ au capital lui-même, surestimant la contribution de $K$ et sous-estimant celle de $A$. En utilisant $K/Y$ (l'intensité capitalistique), on isole la contribution \guillemotleft~pure~\guillemotright\ du capital.
\end{reponse}

\qcm{Selon la comptabilité du développement, le principal facteur expliquant les écarts de revenu entre pays riches et pays pauvres est :}
\choix{
   \item Le stock de capital physique ($K$)
   \item Le taux de participation au marché du travail ($L/N$)
   \item La productivité totale des facteurs ($A$)
   \item La taille de la population ($N$)
}

\begin{reponse}
\textbf{(c)} Les exercices de comptabilité du développement montrent que la PTF ($A$) explique environ \textbf{50\,\%} de la variation du PIB par habitant entre pays. Le capital physique et le taux d'emploi expliquent le reste, mais leur contribution est bien moindre. Les pays pauvres ne sont pas pauvres principalement parce qu'ils manquent de machines, mais parce que leurs institutions, technologies et organisation de la production sont moins efficaces.
\end{reponse}

\qcm{Selon le modèle de Solow, un pays qui augmente durablement son taux d'épargne (et donc son taux d'investissement) :}
\choix{
   \item Connaît une croissance permanente plus rapide du PIB par habitant
   \item Atteint un niveau de PIB par habitant plus élevé, mais la croissance finit par s'arrêter
   \item Voit son PIB par habitant diminuer à cause de la dépréciation
   \item N'a aucun effet sur le PIB par habitant
}

\begin{reponse}
\textbf{(b)} Dans le modèle de Solow, une hausse du taux d'épargne augmente l'investissement, ce qui fait croître le stock de capital. Mais les \textbf{rendements décroissants} du capital impliquent que l'investissement supplémentaire finit par être absorbé par la dépréciation. L'économie converge vers un nouvel état stationnaire \textbf{plus élevé}, mais la croissance s'arrête une fois cet état atteint. Seule la productivité ($A$) peut soutenir la croissance de long terme.
\end{reponse}

\qcm{Dans le modèle de Solow, une augmentation de la croissance démographique entraîne :}
\choix{
   \item Un PIB par habitant plus élevé à l'état stationnaire
   \item Un PIB par habitant plus faible à l'état stationnaire
   \item Aucun effet sur le PIB par habitant
   \item Un taux de croissance permanemment plus élevé
}

\begin{reponse}
\textbf{(b)} Dans l'analogie de la baignoire, une croissance démographique plus rapide augmente le \textbf{débit du drain} : il faut plus d'investissement juste pour maintenir le capital par travailleur ($K/L$) constant face à l'arrivée de nouveaux travailleurs. Le résultat est un état stationnaire avec un $K/L$ et un $Y/N$ plus \textbf{faibles}.
\end{reponse}

\qcm{Dans l'analogie de la baignoire du modèle de Solow, le robinet représente l'investissement. Le drain représente :}
\choix{
   \item La consommation des ménages
   \item Les exportations nettes
   \item La dépréciation du capital et la croissance de la population
   \item Les impôts prélevés par le gouvernement
}

\begin{reponse}
\textbf{(c)} Le drain a deux composantes : (1) la \textbf{dépréciation} ($\delta$), car les machines s'usent et doivent être remplacées, et (2) la \textbf{croissance de la population} ($n$), car de nouveaux travailleurs arrivent et il faut équiper chacun en capital pour maintenir $K/L$ constant. Le débit du drain est proportionnel au stock de capital : $(\delta + n) \times K/L$.
\end{reponse}

\qcm{L'argument de \guillemotleft~réplication~\guillemotright\ (construire un McDonald identique de l'autre côté de la rue) illustre la propriété :}
\choix{
   \item Des rendements décroissants du capital
   \item Des rendements d'échelle constants
   \item Des rendements d'échelle croissants
   \item De la productivité totale des facteurs
}

\begin{reponse}
\textbf{(b)} L'argument de réplication montre que si on double \textbf{tous} les intrants ($K$ et $L$), on peut exactement reproduire l'usine existante et donc doubler la production. C'est la définition des rendements d'échelle constants : $F(2K, 2L) = 2F(K, L)$. Attention : cela ne contredit pas les rendements décroissants, qui s'appliquent quand on augmente un \textbf{seul} facteur.
\end{reponse}

\qcm{La convergence économique (les pays pauvres rattrapent les pays riches) est observée :}
\choix{
   \item Entre tous les pays du monde sans exception
   \item Uniquement entre les pays de l'OCDE et les États américains
   \item Entre les pays qui partagent des fondamentaux similaires (institutions, éducation, etc.)
   \item Uniquement en Asie
}

\begin{reponse}
\textbf{(c)} La convergence \textbf{conditionnelle} est bien documentée : parmi les pays qui partagent des fondamentaux similaires (institutions, éducation, ouverture commerciale), les plus pauvres croissent plus vite. C'est le cas entre les États américains, entre les pays de l'OCDE, et entre les pays asiatiques. Cependant, la convergence \textbf{absolue} (entre tous les pays) n'est pas observée : certains pays pauvres stagnent, notamment en raison d'institutions faibles.
\end{reponse}

\qcm{Après la Deuxième Guerre mondiale, l'Allemagne et le Japon ont connu une croissance très rapide. Selon le modèle de Solow, cela s'explique par le fait que :}
\choix{
   \item La guerre a amélioré leur productivité ($A$)
   \item La destruction de capital a placé ces pays bien en dessous de leur état stationnaire
   \item Leur taux d'épargne a augmenté de façon permanente après la guerre
   \item Leur population a fortement diminué pendant la guerre
}

\begin{reponse}
\textbf{(b)} La guerre a massivement détruit le stock de capital ($K/L$ a chuté), mais les \textbf{fondamentaux} (institutions, éducation, productivité) n'ont pas fondamentalement changé. Loin de l'état stationnaire, le rendement marginal du capital est élevé : dans la baignoire, le robinet coule beaucoup plus vite que le drain. Le résultat est une phase de \textbf{rattrapage} rapide, jusqu'à ce que l'économie retrouve son état stationnaire d'avant-guerre.
\end{reponse}

\qcm{Malgré l'accélération de l'innovation technologique (internet, IA, etc.), la croissance de la PTF dans les économies avancées a :}
\choix{
   \item Fortement accéléré depuis les années 1970
   \item Ralenti depuis les années 1970
   \item Été exactement stable depuis 50 ans
   \item Accéléré en Europe mais ralenti en Amérique du Nord
}

\begin{reponse}
\textbf{(b)} C'est le \guillemotleft~paradoxe de la productivité~\guillemotright\ : \guillemotleft~on voit des ordinateurs partout, sauf dans les statistiques de productivité~\guillemotright. Malgré des avancées technologiques spectaculaires, la croissance de la PTF a \textbf{ralenti} dans la plupart des économies avancées depuis les années 1970. Le Canada est particulièrement touché par ce phénomène.
\end{reponse}

\qcm{Dans l'identité de Kaya, lequel des quatre leviers a le plus contribué à la réduction des émissions dans les économies avancées?}
\choix{
   \item La réduction de la population ($N$)
   \item La baisse du PIB par habitant ($Y/N$)
   \item La réduction de l'intensité énergétique ($E/Y$)
   \item La réduction de l'intensité carbone ($\text{CO}_2/E$)
}

\begin{reponse}
\textbf{(c)} L'identité de Kaya décompose les émissions : $\text{CO}_2 = N \times (Y/N) \times (E/Y) \times (\text{CO}_2/E)$. Dans les pays avancés, la population augmente lentement et le PIB par habitant continue de croître. Le levier principal a historiquement été la baisse de l'\textbf{intensité énergétique} ($E/Y$) : produire la même valeur avec moins d'énergie grâce à l'efficacité énergétique, la désindustrialisation et la tertiarisation. Plus récemment, la baisse de l'intensité carbone ($\text{CO}_2/E$) via les énergies renouvelables joue un rôle croissant.
\end{reponse}

\qcm{Selon le paradoxe de Jevons (effet rebond), une amélioration de l'efficacité énergétique :}
\choix{
   \item Réduit toujours la consommation totale d'énergie
   \item Peut augmenter la consommation totale d'énergie si la demande est suffisamment élastique
   \item N'a aucun effet sur la consommation d'énergie
   \item Réduit la consommation d'énergie uniquement dans les pays en développement
}

\begin{reponse}
\textbf{(b)} Le paradoxe de Jevons montre qu'une meilleure efficacité réduit le \textbf{coût d'usage} d'une ressource, ce qui peut stimuler la demande au point d'\textbf{annuler} les gains d'efficience. L'exemple classique est l'éclairage LED : chaque ampoule consomme 80\,\% moins d'énergie, mais nous éclairons tellement \textbf{plus} de choses qu'avant que la consommation totale a peu diminué. L'ampleur de l'effet rebond dépend de l'élasticité-prix de la demande.
\end{reponse}

% ============================================================================
\section{Vrai, faux ou incertain}
% ============================================================================

\setcounter{shortcounter}{0}

\vfi{Un pays peut soutenir une croissance permanente de son PIB par habitant en investissant massivement dans le capital physique.}

\begin{reponse}
\textbf{Faux.} Les rendements décroissants du capital font que chaque unité de capital supplémentaire ajoute de moins en moins à la production. Dans le modèle de Solow, l'investissement finit toujours par être absorbé par la dépréciation : l'économie converge vers un état stationnaire. Seule la \textbf{productivité} ($A$) --- technologie, institutions, organisation --- peut soutenir la croissance de long terme.
\end{reponse}

\vfi{Si la productivité totale des facteurs ($A$) augmente de 1\,\%, le PIB par habitant augmente d'environ 1\,\%.}

\begin{reponse}
\textbf{Faux.} Le PIB par habitant augmente de $\frac{1}{1-\alpha} \approx 1.43\,\%$ (avec $\alpha = 0.3$). Le coefficient $\frac{1}{1-\alpha}$ est un \textbf{amplificateur} : quand $A$ augmente, la production augmente, ce qui stimule l'investissement, ce qui augmente le capital, ce qui augmente encore la production. Le cycle $A\uparrow \to Y\uparrow \to K\uparrow \to Y\uparrow$ amplifie l'effet initial.
\end{reponse}

\vfi{Un pays plus grand (plus de population) a nécessairement un niveau de vie plus élevé.}

\begin{reponse}
\textbf{Faux.} Les rendements d'échelle constants impliquent que la \textbf{taille} d'un pays ne détermine pas son \textbf{niveau de vie}. Si on double $K$, $L$ et $N$ simultanément, $Y$ double, mais $Y/N$ reste inchangé. Des pays très petits comme le Luxembourg, Singapour ou le Liechtenstein comptent parmi les plus riches du monde, tandis que des pays très peuplés comme l'Inde ou le Nigeria ont un PIB par habitant faible.
\end{reponse}

\vfi{Puisque la PTF est le principal moteur de la croissance, un pays devrait investir uniquement en R\&D et non en capital physique.}

\begin{reponse}
\textbf{Faux.} Bien que la PTF soit le seul moteur de la croissance \textbf{soutenue} à long terme, le capital physique détermine le \textbf{niveau} de vie. Un pays sans capital suffisant sera pauvre, même avec une bonne productivité. De plus, le capital \textbf{amplifie} les gains de productivité : quand $A$ augmente de 1\,\%, le PIB par habitant augmente de $\frac{1}{1-\alpha} \approx 1.43\,\%$ grâce au cycle $A \uparrow \to Y \uparrow \to K \uparrow \to Y \uparrow$. L'investissement en capital est complémentaire à la R\&D, pas un substitut.
\end{reponse}

\vfi{Si une guerre détruit une grande partie du capital d'un pays, le modèle de Solow prédit que le PIB par habitant sera \textbf{permanemment} plus bas qu'avant la guerre.}

\begin{reponse}
\textbf{Faux.} La guerre détruit le stock de capital ($K/L$ chute), mais si les \textbf{fondamentaux} du pays (institutions, productivité, taux d'épargne) restent intacts, l'état stationnaire n'a pas changé. Loin de l'état stationnaire, le rendement marginal du capital est élevé, ce qui stimule l'investissement. L'économie traverse une phase de \textbf{rattrapage} rapide et revient à son état stationnaire d'avant-guerre. C'est exactement ce qu'on a observé en Allemagne et au Japon après 1945.
\end{reponse}

\vfi{Tous les pays pauvres finiront par rattraper les pays riches.}

\begin{reponse}
\textbf{Faux.} Le modèle de Solow prédit la convergence uniquement entre pays qui partagent les mêmes \textbf{fondamentaux} (institutions, éducation, taux d'épargne). C'est la convergence \textbf{conditionnelle}. On l'observe entre les États américains et au sein de l'OCDE, mais \textbf{pas} entre tous les pays du monde. Plusieurs pays pauvres stagnent depuis des décennies en raison d'institutions faibles, de conflits ou de mauvaise gouvernance.
\end{reponse}

\vfi{Un pays qui épargne davantage qu'un autre aura un \textbf{taux de croissance} du PIB par habitant plus élevé à long terme.}

\begin{reponse}
\textbf{Faux.} Il faut distinguer \textbf{niveau} et \textbf{taux de croissance}. Un taux d'épargne plus élevé mène à un état stationnaire avec un capital par travailleur ($K/L$) plus élevé, donc un \textbf{niveau} de PIB par habitant plus élevé. Mais une fois l'état stationnaire atteint, le \textbf{taux de croissance} du PIB par habitant ne dépend que de la croissance de la productivité ($\pctchg{A}$), pas du taux d'épargne.
\end{reponse}

\vfi{La PTF peut être directement mesurée à partir des comptes nationaux.}

\begin{reponse}
\textbf{Faux.} La PTF est calculée comme un \textbf{résidu} : $\pctchg{A} = \pctchg{Y} - \alpha \cdot \pctchg{K} - (1-\alpha) \cdot \pctchg{L}$. On mesure d'abord $Y$, $K$, $L$ et $\alpha$, puis on déduit $A$ par soustraction. C'est une \guillemotleft~mesure de notre ignorance~\guillemotright\ : $A$ capte \textbf{tout} ce qui affecte la production au-delà du capital et du travail (technologie, institutions, savoir-faire, organisation).
\end{reponse}

\vfi{Les améliorations d'efficacité énergétique suffisent à elles seules pour réduire les émissions de CO$_2$ à long terme.}

\begin{reponse}
\textbf{Incertain.} L'efficacité énergétique réduit l'intensité énergétique ($E/Y$), ce qui tend à réduire les émissions. Mais le \textbf{paradoxe de Jevons} (effet rebond) montre que la baisse du coût d'usage peut stimuler la demande, annulant une partie des gains. Par exemple, les ampoules LED consomment 80\,\% moins, mais on éclaire beaucoup plus de choses. Pour garantir une baisse des émissions, il faut compléter l'efficacité par des politiques de prix (taxe carbone) ou des réglementations.
\end{reponse}

\vfi{Plusieurs pays européens ont réussi à réduire leurs émissions de CO$_2$ tout en maintenant une croissance économique positive.}

\begin{reponse}
\textbf{Vrai.} C'est le phénomène de \textbf{découplage} entre croissance économique et émissions. Grâce à la baisse de l'intensité énergétique ($E/Y$) et de l'intensité carbone ($\text{CO}_2/E$), plusieurs pays européens ont réduit leurs émissions tout en augmentant leur PIB. Dans l'identité de Kaya, $Y/N$ a augmenté mais $E/Y$ et $\text{CO}_2/E$ ont baissé suffisamment pour que les émissions totales diminuent.
\end{reponse}

\vfi{La croissance du miracle asiatique (Corée du Sud, Taïwan, etc.) entre 1970 et 1990 était principalement tirée par des gains de productivité ($A$).}

\begin{reponse}
\textbf{Faux.} La décomposition de la croissance montre que le miracle asiatique était largement tiré par l'\textbf{accumulation de capital} (taux d'investissement de 30--40\,\% du PIB), avec une contribution modeste de la PTF. C'est un modèle de croissance par \guillemotleft~transpiration~\guillemotright\ plutôt que par \guillemotleft~inspiration~\guillemotright. Le modèle de Solow prédit que cette croissance devait finir par ralentir --- c'est exactement ce qu'on observe depuis les années 2000.
\end{reponse}

\vfi{Les pays riches en ressources naturelles connaissent systématiquement une croissance plus rapide que les autres.}

\begin{reponse}
\textbf{Faux.} C'est le paradoxe de la \guillemotleft~malédiction des ressources~\guillemotright. Plusieurs mécanismes expliquent pourquoi les ressources naturelles peuvent \textbf{freiner} la croissance : la maladie hollandaise (appréciation du taux de change qui nuit au secteur manufacturier), les institutions faibles (la rente encourage la corruption et le \textit{rent-seeking}), et la volatilité des prix des matières premières. Des contre-exemples existent (Norvège, Canada), mais ce sont des pays avec des \textbf{institutions fortes}.
\end{reponse}

\end{document}
